\section{Realization by a bimodule complex}
\label{sec:realization}

We retain the algebra \(B\), its relation basis \(\mathcal R\), and the
module \(M(Y)\) from the preceding section. The fixed lifted diagram
will now give a single tensor functor realizing the selection process.

\begin{theorem}\label{thm:realization}
Let \(k\) be a field, let \(R\) be a finitely presented \(k\)-algebra,
let \(e\in R\) be a central idempotent, and let
\(\alpha:R\to R\) be a unital endomorphism.
Use the finite-dimensional algebra \(B\) and the modules \(M(Y)\)
constructed in the preceding section, and put
\(H(Y)={}_{\alpha}(eY)\).
There is a single bounded complex \(P\) of finite-dimensional
\(B\)-bimodules, termwise projective on the right, such that for
every finite-dimensional left \(R\)-module \(Y\) one has
\begin{equation}\label{eq:realization}
 P\otL_B M(Y)
 \simeq M(H(Y))\oplus M(H(Y))[3]
 \qquad\text{in }\D^b(B\text{-}\modu).
\end{equation}
All choices defining \(P\) are made independently of \(Y\).
\end{theorem}

To prove the theorem, we first replace the diagram of
\Cref{prop:lifted-diagram} by a complex of \(B\)-bimodules.
Homotopy actions of algebras do not admit such replacements in general:
higher coherence conditions can obstruct them; see
\cite[Sections~2 and~6]{Keller2000}.  Here the three-vertex shape allows
us to verify the construction directly. Keller's homotopy-corrected
complexes and bar construction \cite[Sections~3 and~7]{Keller2000}
provide the general methodological context; the finite construction
below uses only the displayed relation homotopies.

\begin{lemma}\label{lem:rectification}
Let \(D_0,D_1,D_2\) be bounded complexes of finitely generated projective
right \(B\)-modules.  Suppose that each arrow \(a:i\to j\) of \(B\)
is assigned a chain map \(f_a:D_i\to D_j\), and that each relation
\[
 \rho=\sum_{b,a}\rho_{ba}\,ba\in\mathcal R
\]
is assigned a right \(B\)-linear map \(h_\rho:D_0\to D_2\) of degree
\(-1\) satisfying
\[
 d_{D_2}h_\rho+h_\rho d_{D_0}
 = f_\rho,\qquad
 f_\rho:=\sum_{b,a}\rho_{ba}f_bf_a .
\]
There is a bounded complex \(P\) of finite-dimensional \(B\)-bimodules,
termwise projective as a right \(B\)-module, together with homotopy
equivalences of right-module complexes
\[
 \iota_i:D_i\longrightarrow \mathbf e_iP \qquad (i=0,1,2),
\]
such that, for every arrow \(a:i\to j\), the two maps
\(a\iota_i\) and \(\iota_jf_a\) from \(D_i\) to \(\mathbf e_jP\)
are chain homotopic.
\end{lemma}

\begin{proof}
We construct the differential, compare its vertex complexes with the
given ones, and finally compare the arrow actions. All tensor products
in this construction are over \(k\). Form the
complexes of bimodules
\begin{align*}
 L'_0&=\bigoplus_{i=0}^2 B\mathbf e_i\otimes_k D_i,\\
 L'_1&=\bigoplus_{a:i\to j}B\mathbf e_j\otimes_k D_i,\\
 L'_2&=\bigoplus_{\rho\in\mathcal R}B\mathbf e_2\otimes_k D_0.
\end{align*}
The left \(B\)-action is on the first tensor factor and the right
\(B\)-action is on the second.  Let
\(d_0,d_1,d_2\) denote the respective internal differentials of these
three complexes.

On the summand of \(L'_1\) indexed by \(a:i\to j\), define
\[
 p(c\otimes v)
  =(ca\otimes v)_i-(c\otimes f_a(v))_j .
\]
Here subscripts indicate summands of \(L'_0\).  On the summand of
\(L'_2\) indexed by \(\rho\), define
\[
 q(c\otimes v)
  =\sum_{b,a}\rho_{ba}
       \bigl((cb\otimes v)_a+(c\otimes f_a(v))_b\bigr).
\]
These are bimodule chain maps of degree zero.  Define also the
bimodule map \(h:L'_2\to L'_0\) of degree \(-1\) by
\[
 h(c\otimes v)=(c\otimes h_\rho(v))_2
\]
on the relation summand \(\rho\).  Expanding the composite \(pq\)
gives
\begin{align*}
 pq(c\otimes v)
 &=\sum_{b,a}\rho_{ba}
   \bigl(
     (cba\otimes v)_0-(cb\otimes f_a(v))_1\\
 &\hspace{40mm}
     +(cb\otimes f_a(v))_1-(c\otimes f_bf_a(v))_2
   \bigr)\\
 &=-(c\otimes f_\rho(v))_2.
\end{align*}
The last equality uses \(\rho=0\) in \(B\).

Set
\[
 P=L'_0\oplus L'_1[1]\oplus L'_2[2]
\]
as a graded bimodule, and equip it with the differential
\[
 d_P=
 \begin{pmatrix}
   d_0 & p    & h\\
   0   & -d_1 & q\\
   0   & 0    & d_2
 \end{pmatrix}.
\]
For example, \(P^n=(L'_0)^n\oplus(L'_1)^{n+1}
\oplus(L'_2)^{n+2}\), so \(p,q,h\) each have total degree one
in this matrix.  The diagonal signs agree with our shift convention.
Since \(p\) and \(q\) commute with the internal differentials,
the \((0,1)\) and \((1,2)\) entries of \(d_P^2\) vanish.
Its remaining off-diagonal entry is
\[
 d_0h+pq+hd_2=0,
\]
by the displayed calculation and the identity for \(h_\rho\).
The diagonal entries vanish as well.  Thus \(P\) is a complex of
bimodules.

Every term \(B\mathbf e_j\otimes_k D_i^n\), viewed as a right
\(B\)-module, is a finite direct sum of copies of \(D_i^n\).
Consequently all terms of \(P\) are finite dimensional and projective
on the right.  There are finitely many summands, and the complexes
\(D_i\) are bounded; hence \(P\) is bounded.

It remains to recover the given vertex complexes and arrow maps from
this strict bimodule complex.

Define
\[
 \iota_i(v)=(\mathbf e_i\otimes v)_i
       \in \mathbf e_iL'_0\subseteq \mathbf e_iP .
\]
This is a chain map and is split injective as a map of graded right
modules.  We prove that its quotient is acyclic.

Call \(j\) the \emph{source index} of a summand whose second tensor
factor is \(D_j\), in any of the three columns.  Give
\(\mathbf e_iP\) the decreasing filtration
\[
 F^r(\mathbf e_iP)
  =\text{the sum of its summands with source index at least }r
  \qquad (0\leq r\leq 3).
\]
These are subcomplexes.  Indeed, the internal differentials and the
path-multiplication parts of \(p,q\) preserve the source index;
the \(f_a\) terms increase it; and \(h\) increases it from \(0\)
to \(2\).  In particular \(F^0=\mathbf e_iP\) and \(F^3=0\).
The quotient by \(\iota_i(D_i)\) inherits this finite filtration.

On the associated graded, only the internal differentials and the
path-multiplication parts remain.  The source-\(i\) part consists
solely of \(\iota_i(D_i)\), since the only path from \(i\) to itself
is the trivial path.  It disappears in the quotient.  Source indices
larger than \(i\) give zero.  If \(j<i\) are adjacent vertices, the
source-\(j\) part is the signed total complex obtained by tensoring
\(D_j\) with the exact sequence
\[
 0\longrightarrow
 k^{\{a:j\to i\}}
 \xrightarrow{\ a\mapsto a\ }
 \mathbf e_iB\mathbf e_j
 \longrightarrow 0.
\]
There are no relations of length one, so the displayed map is an
isomorphism.

For the remaining pair \((j,i)=(0,2)\), that part is the signed total
complex of \(D_0\) tensored with
\[
 0\longrightarrow k^{\mathcal R}
 \longrightarrow k^{\{(b,a):a:0\to1,\ b:1\to2\}}
 \longrightarrow \mathbf e_2B\mathbf e_0
 \longrightarrow 0.
\]
The three nonzero horizontal terms have degrees \(-2,-1,0\).
The first map sends the basis vector indexed by \(\rho\) to
\(\sum_{b,a}\rho_{ba}(b,a)\); the second sends \((b,a)\) to \(ba\).
The first map is injective because \(\mathcal R\) is a basis of the
relation space.  Its image is the kernel of the second map: all
defining relations have length two, and no path of positive length
can enter vertex \(0\) or leave vertex \(2\).  Thus the ideal they
generate, in this path space, is exactly their linear span.  The
second map is surjective because length-two paths span
\(\mathbf e_2B\mathbf e_0\).

These exact complexes of vector spaces are split, and tensoring
with a bounded complex \(D_j\) gives an acyclic total complex.
All associated graded pieces of the quotient are therefore acyclic.
Induction through the finite filtration proves that the quotient
itself is acyclic.  Thus \(\iota_i\) is a quasi-isomorphism.
Its cone is a bounded acyclic complex of projective right modules,
and is contractible: starting in the largest nonzero degree, split
the surjection onto its projective last term and continue downward.
It follows that \(\iota_i\) is a homotopy equivalence.

Finally, for \(a:i\to j\) define a right-linear map of degree \(-1\)
\[
 H_a:D_i\longrightarrow\mathbf e_jP,\qquad
 H_a(v)=(\mathbf e_j\otimes v)_a
          \in\mathbf e_jL'_1[1].
\]
The internal differential on this summand is \(-d_{D_i}\);
its contribution cancels that of \(H_ad_{D_i}\).  The remaining
component is \(p\), giving
\[
 d_PH_a+H_ad_{D_i}
   =a\iota_i-\iota_jf_a .
\]
This is the required arrow homotopy.
\end{proof}

\begin{remark}
The proof also isolates the reason that higher relation homotopies
are unnecessary here.  Every relation runs from vertex \(0\) to
vertex \(2\), and cannot be composed with another nonidentity path.
The three-column differential has only the single square condition
involving \(h\) that was checked above.  There is no further column
or relation overlap that could introduce another coherence equation.
\end{remark}

\begin{proof}[Proof of \Cref{thm:realization}]
Choose the fixed complexes \(D_i\), arrow maps \(f_a\), and relation
homotopies \(h_\rho\) of \Cref{prop:lifted-diagram}, and apply
\Cref{lem:rectification}.  Since \(P\) and the \(D_i\) are bounded
complexes of projective right modules, their ordinary tensor
products with \(M(Y)\) compute the derived tensor products.
Moreover, tensoring preserves the homotopy equivalences \(\iota_i\)
and their arrow homotopies.  Thus the vertex diagram of
\[
 T_Y:=P\otL_B M(Y)
\]
in the derived category of \(k\)-vector spaces is isomorphic to the
evaluated diagram of \Cref{prop:lifted-diagram}.

At each vertex this diagram is \(eY\oplus eY[3]\).  The arrows
\(s,s'\) act by identity maps, and \(a_h,a'_h\) act by
\(\alpha(x_h)\) on each summand.  Applying cohomology, with its
induced arrow maps, therefore gives isomorphisms of left
\(B\)-modules
\[
 H^n(T_Y)\cong
 \begin{cases}
  M(H(Y)),& n=0,-3,\\
  0,& n\ne0,-3.
 \end{cases}
\]
In particular these identifications specify the full \(B\)-action,
not just the vector spaces at the vertices.

Write \(A_{-3}=H^{-3}(T_Y)\) and \(A_0=H^0(T_Y)\).
The canonical truncation triangle is
\[
 A_{-3}[3]\longrightarrow T_Y\longrightarrow A_0
 \xrightarrow{\ \delta\ } A_{-3}[4].
\]
The connecting morphism belongs to
\[
 \Hom_{\D^b(B\text{-}\modu)}(A_0,A_{-3}[4])
   =\Ext_B^4(A_0,A_{-3})=0,
\]
where the vanishing follows from the bound
\(\gldim B\leq2\) in \Cref{lem:directed-bound}.
The triangle consequently splits, giving
\(T_Y\simeq A_0\oplus A_{-3}[3]\), which is
\eqref{eq:realization}.
The complexes, chain maps, and homotopies were fixed before choosing
\(Y\), as was the rectification constructed from them.
\end{proof}

\begin{corollary}
Put \(\Phi=P\otL_B-\).  For every integer \(j\geq0\) and every
finite-dimensional left \(R\)-module \(Y\), there is an isomorphism
\[
 \Phi^j(M(Y))
 \simeq
 \bigoplus_{q=0}^j
   \bigl(M(H^{\circ j}(Y))[3q]\bigr)^{\oplus\binom{j}{q}} .
\]
\end{corollary}

\begin{proof}
The case \(j=0\) is immediate.  For the induction step, apply
\(\Phi\) to an isomorphism for \(j\); derived tensor preserves
isomorphisms, finite direct sums, and shifts.  Apply
\eqref{eq:realization} with \(H^{\circ j}(Y)\) in place of \(Y\) to each
summand, and combine the two resulting multiplicities using
Pascal's identity.  Each \(H^{\circ j}(Y)\) is finite dimensional.
This argument uses the isomorphisms only one object at a time;
no natural choice of the truncation splittings is required.
\end{proof}
